$\Cl(1,n{-}1)$ is the spacelike extension of $\Cl(1,n{-}2)$, an even-dimensional algebra with $p - q = 3 - n$, which is $\equiv 2 \pmod 4$ iff $n \equiv 1 \pmod 4$ and $\equiv 0$ iff $n \equiv 3$; Theorem \claimref{rule2.where-obstruction-sits-commuting-square}(a) gives the first claim and the complex type of the other case. $\Cl(n{-}1,1)$ is the extension of $\Cl(n{-}2,1)$, $p - q = n - 3$, by a generator of square $+1$, obstructed iff $n - 3 \equiv 0$, i.e.\ $n \equiv 3 \pmod 4$, complex type otherwise. For the complexification: by Theorem \claimref{std-algebra.complexification-is-the-odd-extension} applied to the even-dimensional $V \oplus \R f$, $\Cl(\hat V)$ is its complexification with central unit $I_{\hat V}$, and the subalgebra generated by $V$ and $I_{\hat V}$ is $\Cl(V) \oplus \Cl(V)I_{\hat V}$ (the summands meet in $0$, since the blades of the second contain $f$), the complexification $\Cl(p,q) \otimes \C$ with complex unit $I_{\hat V}$; in it the substitution $e_k \mapsto I_{\hat V}e_k$ on the generators of square $-1$ turns the relations of $\Cl(p,q)$ into those of $\Cl(n,0)$, so $\Cl(p,q) \otimes \C$ depends only on $n$ (\claimref{std-algebra.complexification}); the volume element remains central, and if $\omega^2 = -1$ then $(I_{\hat V}\omega)^2 = +1$, so $\omega'$ as defined is a central involution and $\tfrac12(1 \pm \omega')$ are orthogonal central idempotents summing to $1$. The script checks the alternation, with its complex-type complement, in dimensions $3, 5, 7, 9, 11$.
